% ============================================================================
% Chapter 13 --- Cryptographic services
% ----------------------------------------------------------------------------
% Purpose : PKI, secrets management, HSM, ACME. Belongs to the Subject and
%           Trust plane (D5.1). Addresses the live-path HSM integration
%           anticipated by D10.6 (HYOK for higher-classification storage).
% Page target: ~3 pages.
% Dependencies: Ch. 5 (PKCS#11, ACME, KMIP, OpenID Federation),
%               Ch. 6 (PKI underpins identity assertions),
%               Ch. 10 (HYOK live-path for storage),
%               Ch. 11 (device certificates and secrets).
% Mirrors layer chapter template established in Ch. 6 (D6.5).
% ----------------------------------------------------------------------------
% Decisions registered in _coherence-EN.md:
%   D13.1 minimum compliance set for cryptographic services
%   D13.2 reference instantiation = step-ca + Vault + SoftHSM/Nitrokey HSM
%   D13.3 alternatives = AD CS + commercial HSM; external CA via ACME
%   D13.4 HYOK live-path requirement from D10.6 addressed here
%   D13.5 cryptographic services sovereignty scores
% ============================================================================

\chapter{Cryptographic services}
\label{ch:crypto}

\begin{caveatbox}[Dependencies]
\noindent This chapter applies the layer template of
\trchap{ch:identity}{} to the cryptographic-services layer. It
underpins the assertions of \trchap{ch:identity}{}, the encryption
at rest of \trchap{ch:storage}{} (including the HYOK live-path of
D10.6), and the device certificates of \trchap{ch:endpoint}{}. The
layer belongs to the \emph{Subject and trust plane}.
\end{caveatbox}

The cryptographic-services layer is the foundation on which the
trustworthy assertions of every other layer rest: the certificates
that authenticate institutional services, the secrets that
applications consume, the keys that protect data at rest. It is
also the layer where errors are the most consequential: a
compromised intermediate certificate authority, an unrotatable key,
or a secrets store whose backup is missing each translate into
incidents whose remediation is institution-wide.

%-----------------------------------------------------------------------------
\section{Functional role and interface contract}
\label{sec:cr-contract}

\begin{contractbox}[Cryptographic services]
\noindent A compliant cryptographic-services layer must:

\begin{itemize}[leftmargin=1.5em,nosep]
  \item operate an institutional certificate authority hierarchy
    whose trust anchor is held in hardware (HSM accessed through
    PKCS\#11~\cite{pkcs11-oasis}) under institutional custody; the
    key-protection policy declares at procurement whether
    trust-anchor keys are wrap-exportable for migration and, where
    they are non-exportable (a common HSM best practice), defines
    the alternative exit artefact: a planned key-rotation and
    re-anchoring rehearsal under dual control;
  \item issue end-entity certificates through
    ACME~\cite{rfc-8555-acme} for machine-readable workflows, and
    through a documented manual or automated procedure for endpoint
    certificates;
  \item provide a secrets-management facility (key-value, dynamic
    secrets, transit encryption) accessible through documented APIs
    and integrable with the policy layer for access decisions;
  \item support a Bring-Your-Own-Key arrangement for any external
    service that consumes institutional cryptographic material;
  \item provide a Hold-Your-Own-Key live-path arrangement for the
    higher-classification data of \trchap{ch:storage}{} (D10.6),
    in which keys never leave the institutional HSM in cleartext;
  \item produce audit events conformant to the canonical TAC
    schema of \trchap{ch:observability}{} (\S\ref{sec:ob-schema}).
\end{itemize}
\end{contractbox}

%-----------------------------------------------------------------------------
\section{Reference instantiation: open-source ACME CA + secrets-management platform + software or hardware HSM}
\label{sec:cr-reference}

\begin{instantiationbox}[Cryptographic services]
\noindent\textbf{Topology.} An open-source ACME-capable certificate
authority --- for example \texttt{step-ca}~\cite{step-ca}
(Smallstep) --- operates the institutional certificate authority,
exposing ACME for machine issuance and standard interfaces for
human-driven workflows. An open-source secrets-management platform
--- for example HashiCorp Vault~\cite{vault} --- provides secrets
management with an envelope-encryption (transit) engine
(consumed by \trchap{ch:storage}{}) and
a PKI engine as a complementary issuance path. The trust anchor of
the CA hierarchy is held in an HSM exposed through PKCS\#11:
representative choices are a software HSM
(such as SoftHSM~\cite{softhsm}) for non-production environments and
a hardware HSM (such as Nitrokey HSM, YubiHSM, or comparable) for
production. For HYOK live-path operation, the storage layer's
encryption requests are forwarded to the HSM through a thin
PKCS\#11 broker that the institution operates.
\end{instantiationbox}

\paragraph{Operational considerations.} The reference instantiation
would typically require a small dedicated staffing
allocation --- a fractional engineering capacity shared with the
security and identity teams, sized to the institution's scale ---
plus institutional governance for the CA hierarchy (root rotation,
intermediate rotation, revocation procedures). The principal
recurring effort is operational discipline rather than ongoing
engineering: most of the value comes from following the
documented procedures, not from continuously evolving the
configuration.

\begin{realworldbox}[Public infrastructure precedents]
\noindent The ACME ecosystem at internet scale, exemplified by
Let's Encrypt and the IETF RFC 8555 specification, demonstrates
that automated certificate issuance is operationally robust. Most
research and education networks operate certificate services for
their members with comparable patterns; specific institutional
configurations should be confirmed against current public
documentation when cited.
\end{realworldbox}

%-----------------------------------------------------------------------------
\section{Alternative~1 --- Active Directory Certificate Services + commercial HSM}
\label{sec:cr-alt1}

\begin{alternativebox}[AD CS]
\noindent The institution operates a Microsoft Active Directory
Certificate Services hierarchy with a commercial HSM (Thales
Luna, Entrust nShield, or comparable) holding the trust anchor.
ACME is exposed through a connector or front-end so that the
certificate-issuance interface to applications matches the
reference contract.
\end{alternativebox}

\paragraph{Where it adds value.} For institutions whose Microsoft
estate already includes AD CS, the operational simplification of
keeping the certificate authority in the same administrative
boundary is substantial. Commercial HSMs have well-documented
certifications (FIPS 140 levels) that may be required by specific
research mandates or regulatory contexts.

\paragraph{Where it constrains the institution.} ADCS expresses
issuance policy in templates whose semantics are
Microsoft-specific; bridging to the standard contract requires the
ACME connector to translate faithfully, and the connector becomes
itself a critical component. Commercial HSMs introduce per-device
licensing and a vendor relationship in the cryptographic-services
layer that the reference instantiation avoids.

\paragraph{When it could be the right choice.} An institution with
mature AD CS operations and either a regulatory or research
requirement that the chosen HSM specifically satisfies, where the
operational continuity outweighs the sovereignty cost on the
technical and financial dimensions.

%-----------------------------------------------------------------------------
\section{Alternative~2 --- External CA via ACME standard}
\label{sec:cr-alt2}

\begin{alternativebox}[External CA via ACME]
\noindent The institution does not operate its own institutional
CA hierarchy beyond what is required for internal services, and
relies on an external commercial CA (Sectigo, DigiCert, GlobalSign)
or on Let's Encrypt for public-facing certificates, all consumed
through the ACME standard. Internal services and HYOK arrangements
remain on the institutional Vault and HSM.
\end{alternativebox}

\paragraph{Where it adds value.} For public-facing certificates, an
external CA whose trust roots are pre-installed in browsers and
operating systems eliminates an institutional concern (browser
trust) that the reference instantiation would have to address
through manual root distribution. Let's Encrypt in particular has
demonstrated that an external CA exposed through ACME is
operationally compatible with the architecture's standards-only
posture.

\paragraph{Where it constrains the institution.} The institution
depends on the external CA's continued operation and on its
acceptance criteria. For public-facing services this dependency is
modest (multiple CAs are interchangeable through ACME); for
research-specific or regulatory contexts it may be more
constraining. The financial dimension is excellent for Let's
Encrypt and predictable for commercial CAs; the technical dimension
inherits the CA's open or proprietary posture.

\paragraph{When it could be the right choice.} As a complement,
not a substitute, for the institutional CA: external CAs handle
public-facing certificates while the institution retains
sovereignty over its internal hierarchy.

%-----------------------------------------------------------------------------
\section{Sovereignty grid applied}
\label{sec:cr-sovereignty}

\begin{table}[H]
\centering\small
\caption{Per-dimension sovereignty profile of the cryptographic-services
instantiations.}
\label{tab:cr-sovereignty}
\begin{tabularx}{\textwidth}{%
  >{\hsize=0.6\hsize\linewidth=\hsize\bfseries\raggedright\arraybackslash}X%
  >{\hsize=1.2\hsize\linewidth=\hsize\centering\arraybackslash}X%
  >{\hsize=1.2\hsize\linewidth=\hsize\centering\arraybackslash}X%
  >{\hsize=1.2\hsize\linewidth=\hsize\centering\arraybackslash}X%
  >{\hsize=0.8\hsize\linewidth=\hsize\centering\arraybackslash}X%
}
\toprule
Dimension & Reference (open-source CA + secrets platform + HSM) & AD CS + commercial HSM & External CA via ACME \\
\midrule
Data          & 3 Robust      & 2 Established & 2 Established \\
Technical     & 3 Robust      & 2 Established & 2 Established \\
Financial     & 2 Established & 1 Partial     & 3 Robust      \\
Operational   & 2 Established & 3 Robust      & 3 Robust      \\
Institutional & 3 Robust      & 2 Established & 2 Established \\
\bottomrule
\end{tabularx}
\end{table}

\noindent The cryptographic-services layer is the layer in which the
sovereignty trade-offs are typically the smallest in absolute terms
(none of the alternatives is structurally below Established on any
dimension when adopted with discipline) and the largest in practical
consequence (a misalignment here propagates to every layer that
depends on cryptographic assertions). The architectural posture is
to keep the trust anchor under institutional custody regardless of
the operational instantiation, and to choose the operational
instantiation according to the institution's existing investments
and regulatory context.

%-----------------------------------------------------------------------------
\section{Regulatory anchoring}
\label{sec:cr-regulatory}

\noindent The cryptographic-services layer is the
trust-anchor and key-custody substrate. It is the principal
anchor of the applicable cybersecurity-baseline regime's
cryptography requirements (in the European Union, NIS2~Art.~21(2)(h)
--- cryptography policies and procedures) and contributes to the
applicable data protection regime's encryption and pseudonymisation
obligations (in the European Union, GDPR~Art.~32(1)(a)).
It is also the operational counterpart of the applicable
electronic-identification and trust-services regime
(in the European Union, eIDAS~2: qualified electronic
signatures, seals and time stamps, when the institution
provides them). The layer
operationalises 27001~A.8.24 (use of cryptography) and
contributes to NIST CSF \textsf{PROTECT} (PR.DS --- data
security, encryption at rest and in transit). Detailed
mapping in \trchap{ch:regulatory-mapping}{}
\S\ref{sec:rm-gdpr}, \S\ref{sec:rm-nis2},
\S\ref{sec:rm-eidas2}, \S\ref{sec:rm-iso},
\S\ref{sec:rm-nist-csf}.

%-----------------------------------------------------------------------------
\section{Common pitfalls}
\label{sec:cr-pitfalls}

\begin{pitfallbox}[Trust anchor without an HSM]
\noindent An institutional CA whose root or intermediate keys live
on a workstation or a server filesystem is one disk image away
from an institution-wide trust collapse. The mitigation is to hold
the trust anchor in an HSM under institutional custody, and to
exercise the recovery procedure periodically as part of
\trchap{ch:conformance}{}.
\end{pitfallbox}

\begin{pitfallbox}[Key escrow drift]
\noindent The institution declares a key-escrow policy (which keys
are escrowed, where, accessible by whom, under which procedure)
and maintains it nominally; the actual escrowed keys, the actual
custodians, and the actual procedure drift from the policy over
years. Discovery typically occurs during a custodian transition
or an incident. The mitigation is to verify the escrow state at
the same cadence as the sovereignty grid (annual + event-driven).
\end{pitfallbox}

\begin{pitfallbox}[ACME outage with no fallback]
\noindent An ACME-driven certificate renewal pipeline whose
upstream CA is temporarily unavailable can produce service
outages: certificates expire because the renewal pipeline cannot
issue new ones. The mitigation is to provision certificates with a
margin against renewal failures (renew well before expiry), to
monitor the renewal pipeline through the observability layer, and
to maintain a documented manual-issuance fallback procedure for
the highest-criticality services.
\end{pitfallbox}

\begin{pitfallbox}[CA hierarchy not documented]
\noindent The institutional CA hierarchy --- root, intermediates,
issuing CAs, their key strengths, their validity periods, their
rotation schedules --- is documented either incompletely or in a
form that is not currently accessible. Replacement or extension
becomes a recurring archaeology project. The mitigation is to
maintain the hierarchy documentation under version control, with
the same review cadence as the policy corpus.
\end{pitfallbox}
